\section{Rigid boundary classes on a torus}
\label{sec:torus-boundary-rigidity}

We now turn to the boundary of the nef cone. A nef class with zero top intersection number may have a smooth semipositive representative and yet leave little freedom to vary it. On a complex torus, the kernel of a constant semipositive form controls this freedom: when the leaves of the kernel are dense, the form is the only closed positive current in its class. The same condition decides whether adding a small fixed K\"ahler form produces metrics whose diameters tend to zero. For K\"ahler forms converging to the boundary in other ways, the diameters can tend to zero even at a class with many semipositive representatives.

Let $X=V/\Lambda$ be a complex torus, where $V\simeq\mathbb C^n$ and $\Lambda$ is a full lattice. Let $\theta$ be a translation-invariant semipositive real $(1,1)$-form of rank $r$ with $1\leq r<n$, and fix a translation-invariant K\"ahler form $\kappa$. The class $\alpha=\{\theta\}_{\mathrm{BC}}$ is nef and has $\int_X\theta^n=0$. It lies on the boundary of the nef cone: a K\"ahler form in $\alpha$ would have positive top intersection number, whereas Stokes' theorem gives every representative the top intersection number $\int_X\theta^n=0$. We call $\alpha$ \emph{rigid} if it contains exactly one closed positive current. Let $K=\ker\theta\subset V$, a complex subspace of dimension $n-r$, and let $H\subset X$ be the closure of its image in $X$. This is a connected real subtorus, which need not be complex. Diameters are computed for the Riemannian metric $g_\eta(v,w)=\eta(v,Jw)$ associated with a positive form $\eta$, where $J$ is the complex structure.

\Needspace{17\baselineskip}
\begin{theorem}\label{thm:flat-torus-radial-rigidity}\index{Nef class!rigid torus boundary}\index{Metric degeneration!dense kernel leaves}
The following are equivalent:
\begin{enumerate}
\item $\alpha$ is rigid;
\item $H=X$;
\item $\operatorname{diam}(X,\theta+t\kappa)\to0$ as $t\downarrow0$.
\end{enumerate}
In the first case the unique current is $\theta$. The third condition is independent of the chosen translation-invariant K\"ahler form $\kappa$.
\end{theorem}
\begin{proof}
\emph{Rigidity when $H=X$.} Suppose first that $H=X$. By \cref{lem:bc-one-one-current-description}, a positive current in $\alpha$ can be written $T=\theta+\ddc u$ with $u$ quasi-plurisubharmonic, in particular $u\in L^1(X)$. For a constant real vector $v\in K$, the coefficient $T(v,Jv)$ is a positive Radon measure. Its total integral against Haar measure is zero, because $\theta(v,Jv)=0$ and the remaining part of the coefficient is a combination of constant-coefficient derivatives of $u$. Hence this coefficient vanishes. Writing $D_v$ for differentiation along $v$ and using that $K$ is complex, we obtain
\[
 (D_v^2+D_{Jv}^2)u=0
 \quad\text{as a distribution for every }v\in K.
\]
For a real linear functional $\ell$ with $\ell(\Lambda)\subset\mathbb Z$, let $\widehat u(\ell)$ be the Fourier coefficient of $u$ at the character $\exp(2\pi\mathrm i\ell)$. The identity above gives
\[
 \bigl(\ell(v)^2+\ell(Jv)^2\bigr)\widehat u(\ell)=0.
\]
Since the image of $K$ is dense, a nonzero character cannot vanish identically on $K$. Thus every nonconstant Fourier coefficient of $u$ is zero, and convolution with a smooth approximate identity shows that $u$ is constant. Hence $T=\theta$.

\emph{Nonrigidity when $H\ne X$.} If $H\ne X$, a nontrivial character of the real torus $X/H$ gives a nonzero $\ell$ that vanishes on $K$. Put $\phi=\cos(2\pi\ell)$. Since $K$ is complex, the form $\ddc\phi$ vanishes whenever one of its arguments lies in $K$. On a fixed complement of $K$, $\theta$ is positive definite and $\ddc\phi$ is uniformly bounded. Consequently, the forms $\theta\pm\varepsilon\ddc\phi$ are semipositive for sufficiently small $\varepsilon>0$. They represent the same class, and they are distinct because a nonconstant smooth function on a compact torus cannot have zero $\ddc$. Hence $\alpha$ is not rigid.

\emph{Collapse when $H=X$.} Suppose again that $H=X$, and fix $\delta>0$. By density and compactness, there are finitely many vectors $k_1,\ldots,k_N\in K$ whose images form a $\delta$-net in the flat metric of $\kappa$. Put $R=\max_j|k_j|_\kappa$. Since $\theta(k_j,Jk_j)=0$, the straight path from $0$ to the image of $k_j$ has $(\theta+t\kappa)$-length $\sqrt t\,|k_j|_\kappa$. For $0<t\leq1$ these metrics are bounded above by a fixed multiple of $\kappa$. Joining this path to a point within $\delta$ therefore gives
\[
 \operatorname{diam}(X,\theta+t\kappa)
       \leq\sqrt t\,R+C\delta,
\]
since by translation invariance the diameter equals the supremum of the distances from $0$. Letting first $t\downarrow0$ for this fixed net and then $\delta\downarrow0$ proves the third condition.

\emph{No collapse when $H\ne X$.} If $H\ne X$, choose $\ell$ as above. Its differential vanishes on $K$, so $|\mathrm d\ell(v)|\leq L|v|_\theta$ for some $L>0$, where $|v|_\theta$ is the seminorm defined by $\theta$. Since $\theta+t\kappa\geq\theta$, the surjective map $\ell\bmod\mathbb Z$ onto the circle of length one is $L$-Lipschitz for every $t>0$. Two points mapping to opposite points of the circle therefore have distance at least $1/(2L)$, and the diameters cannot tend to zero. This completes the proof of the equivalences. Since the second condition does not involve $\kappa$, neither does the third.
\end{proof}

On a two-dimensional square torus, the slope of the kernel line decides rigidity.

\begin{example}\label{ex:irrational-torus-rigidity}\index{Nef class!irrational rigid torus class}
On $X=(\mathbb C/(\mathbb Z+\mathrm i\mathbb Z))^2$, for $a\in\mathbb R$ set
\[
 \theta_a=\frac{\mathrm i}{2}
       (\mathrm dz_2-a\,\mathrm dz_1)
       \wedge(\mathrm d\bar z_2-a\,\mathrm d\bar z_1).
\]
Its class is rigid exactly when $a$ is irrational. In that case its ray is irrational, that is, no nonzero real multiple of its de Rham class is rational, and every family $\theta_a+t\kappa$ as above collapses to a point in diameter. For rational $a$ the class is not rigid, and these diameters have a positive lower bound.
\end{example}
\begin{proof}
The kernel is the complex line $z_2=a z_1$. Write $z_j=x_j+\mathrm i y_j$. An integral character has the form
\[
 \ell=m_1x_1+m_2x_2+n_1y_1+n_2y_2,
 \qquad m_j,n_j\in\mathbb Z,
\]
and it vanishes on the kernel exactly when $m_1+a m_2=n_1+a n_2=0$. For irrational $a$, only the zero character does so, hence the image of the kernel is dense. For rational $a$, a nonzero such character exists. Now apply \cref{thm:flat-torus-radial-rigidity}. The periods of $\theta_a$ on the coordinate two-tori in the $(x_2,y_2)$ and $(x_1,y_2)$ directions are $1$ and $-a$, respectively. If a nonzero multiple of its class were rational, the first period would make that multiple rational, and the second would then make $a$ rational. This proves the assertion about the ray.
\end{proof}

The fixed positive direction $\kappa$ in the theorem cannot be replaced by an arbitrary convergent family of K\"ahler forms, even when all the metrics are flat and Ricci-flat.

\Needspace{16\baselineskip}
\begin{example}\label{ex:unrestricted-torus-collapse}\index{Metric degeneration!unrestricted class convergence}
On the same square torus let $\theta_0=\frac{\mathrm i}{2}\mathrm dz_2\wedge\mathrm d\bar z_2$ and, for $t>0$, put
\[
 \beta_t=\frac{\mathrm i}{2}\left[
    (t\,\mathrm dz_1+\mathrm dz_2)
      \wedge(t\,\mathrm d\bar z_1+\mathrm d\bar z_2)
       +t^2\mathrm dz_2\wedge\mathrm d\bar z_2\right].
\]
These are flat Ricci-flat K\"ahler forms, $\beta_t\to\theta_0$ smoothly, and
\[
 \operatorname{diam}(X,\beta_t)\leq t.
\]
The limiting class $\{\theta_0\}_{\mathrm{BC}}$ is not rigid.
\end{example}
\begin{proof}
The coefficient matrix of $\beta_t$ is $\left(\begin{smallmatrix}t^2&t\\t&1+t^2\end{smallmatrix}\right)$, with determinant $t^4>0$. Its coefficients are constant, which gives flatness and Ricci-flatness. Represent a point by $(z_1,z_2)$ with real and imaginary parts in $[-1/2,1/2]$, and choose a Gaussian integer $k$ nearest to $-z_1-z_2/t$. The displacement $(z_1+k,z_2)$ represents the same point of the torus and has squared metric norm
\[
 |t(z_1+k)+z_2|^2+t^2|z_2|^2
       \leq t^2/2+t^2/2=t^2.
\]
The straight path along this displacement proves the bound for the distance from zero, and translation invariance gives the diameter bound. For $0<|s|<1/\pi$, the form $\theta_0+\ddc(s\cos(2\pi x_2))=(1-\pi s\cos(2\pi x_2))\theta_0$ is another semipositive representative, so the limiting class is not rigid. Finally, the coefficient matrix of $\beta_t-\theta_0$ has eigenvalues $t^2\pm t$, so this difference is indefinite for $0<t<1$: the forms $\beta_t$ do not approach $\theta_0$ from above.
\end{proof}
